\chapter{Search the tool index (menu 5)}
\label{ch:menu5}
\menulabel{5}

\section{Purpose}

Find external tools relevant to f-block and multireference work: what each one
is for, how \fbk{} intends to relate to it, whether it is actually used now,
and its licence.

\section{What you need}

Keywords (several words, space separated; all must match). The same search is
available from the command line: \cmd{search KEYWORDS}.

\section{Walkthrough}

\lstinputlisting[style=fbkcode, firstline=54, lastline=59,
  title={The menu-5 example session}]
  {examples/expected/m5_tools.txt}

\section{What you get}

For every hit: the tier, the status, the licence, the purpose, and the note
with integration caveats. The example entry shows the two independent
dimensions to read:

\begin{center}
\begin{tabularx}{\textwidth}{@{}lX@{}}
\toprule
Dimension & Meaning \\
\midrule
relation (\textbf{tier}) & how we intend to use it: \textbf{A} absorb (file or
  algorithm level), \textbf{B} interface (generate its input, read its output),
  \textbf{C} index (calculation kernels, algorithmic borrowing included) \\
status & whether it is used \emph{now}: \code{active} or
  \code{registered} (``registered'' is a dot on the map, not a feature) \\
\bottomrule
\end{tabularx}
\end{center}

The MOKIT entry is an A-tier tool with status \code{registered}: the project
takes only its ``hand orbitals over'' step and the UNO/GVB window idea, not the
package. Licences the project has not verified are marked ``to verify'' --
never asserted.

\section{How to read the numbers}

\begin{readbox}
Do not read the index as a feature list. ``Registered'' means the tool is on
the map (with a source and a purpose); it does not mean \fbk{} uses it. The
program's own mechanical gate is explicit: the test suite asserts that the
\code{active} set is exactly the one component that is genuinely wired in
(the autoCAS entropy protocol), so any future promotion of a tool to
\code{active} must be accompanied by an actual integration and a test change.
\end{readbox}

\section{Worked example}

Search \code{lanthanide DMET} to find the lanthanide-embedding entries, then
\code{magnetic} for the magnetic-property chain tools. Multi-word searches
require all words to match somewhere in the record.

\section{Caveats, refusals and related sections}

\begin{refusalbox}
\begin{itemize}
  \item The index carries no installation actions and downloads nothing.
  \item A search with no hits says so plainly; an empty keyword is refused.
\end{itemize}
\end{refusalbox}

Related: \menu{6} for the full entry of one tool (onboarding notes).
